\section{Conclusion}\label{sec:Conclusion}

\begin{frame}
    \frametitle{¿Eso es todo?}

    ¿Hasta ahí llegan HTML y CSS?

    No, se puede hacer muchísimo más. Algunas cosas interesantes que se pueden hacer 
    solo con html y css:

    \begin{itemize}
        \item \textbf{Animaciones}: Se puede usar la propiedad \inlinecode{transition} 
              para agregar transiciones cuando cambian las propiedades de los elementos.
              También se puede combinar la propiedad \inlinecode{animation} con la 
              regla \inlinecode{@keyframes} para animar propiedades de forma super 
              específica.
        \item \textbf{Diseño responsive}: Con la regla \inlinecode{@media} se puede 
              definir reglas condicionalmente según el dispositivo en el que se esté 
              mostrando el contenido. Cualquier sitio mínimamente moderno debería 
              usar esta regla para verse adecuadamente en cualquier pantalla.
        \item \textbf{Interactividad}: Si bien esto se suele delegar a JS, HTML
              ahora tiene unas cuantas herramientas como \inlinecode{details} y 
              \inlinecode{dialog} que encapsulan una buena parte del comportamiento.
    \end{itemize}

\end{frame}

\begin{frame}
    \frametitle{Siguientes pasos. Más allá de HTML y CSS.}

    Pero definitivamente no hace falta quedarnos solo con HTML y CSS. El siguiente 
    paso es JavaScript. 

    Si bien antes mencionamos que \emph{técnicamente} se puede hacer cualquier cosa 
    con CSS y HTML, lo normal es usar JavaScript, que sí es un lenguaje de programación 
    de propósito general. 
    
    Aprender JS en realidad es aprender a programar, cosa que les va a tocar hacer a 
    lo largo de la carrera. Las particularidades del lenguaje tienen mucho que ver 
    con que su propósito principal es el de ser el lenguaje de los sitios web. 

    También es muy popular TypeScript. Básicamente es JavaScript pero estáticamente 
    tipado.

    Más recientemente, WebAssembly apareció como una especie de alternativa / 
    extensión de todo esto.

\end{frame}

\begin{frame}
    \frametitle{Herramientas, frameworks!}

    Es probable que hayan escuchado hablar de React, Angular, Tailwind, Bootstrap, o 
    algún otro framework popular. 
    
    La realidad es que hay un montón de cosas que se repiten en muchísimos sitios web.
    Menús, barras de navegación, contenedores, tarjetas, botones, etc. Todo esto 
    puede ser bastante molesto de armar desde 0, por lo que regularmente aparecen 
    herramientas para evitar reinventar la rueda cada vez.

    Hay herramientas bien específicas que resuelven problemas chiquitos, pero las 
    más populares son los frameworks. Estos frameworks muchísimas partes útiles 
    para tomar y reutilizar, haciendo que el desarrollo del sitio web sea mucho más 
    ameno. No solo eso, sino que también vienen con una ``filosofía'' de desarrollo. 
    Cuando se usan ``bien'', todos los sitios hechos con un mismo framework deberían 
    parecerse bastante desde el punto de vista de la implementación.
    
    Este concepto no es exclusivo para el desarrollo web, existe en todos los campos 
    de la programación e incluso por fuera del mismo.

\end{frame}

\begin{frame}
    \frametitle{No nos olvidemos del backend}

    Como mencionamos al principio de la clase, el front solo es una cara del desarrollo web. 
    El back es puede ser más o menos importante dependiendo del caso. 

    En la carrera de computación vemos un montón de temas que conforman la base de lo 
    que después se usa para desarrollar la parte que no se ve de los sitios.

    Un buen lugar para empezar es entender cómo se comunican dos programas entre sí. Para 
    esto se usan APIs, tema del cual tenemos otro taller!

\end{frame}

\begin{frame}
    \frametitle{Hasta acá}

    Esperamos que hayan disfrutado el taller y que se hayan llevado algo útil. También 
    esperamos haberles despertado la curiosidad y haberles dado ganas de profundizar en 
    los temas que de desarrollo web y las cosas que hacemos en la facultad.

    Recuerden que hay una comunidad enorme en esta facultad y en este departamento, por 
    lo que siempre se pueden acercar y preguntar lo que quieran y va a haber gente para 
    contestarles con la mejor de las ondas.

    También pueden conectar con mucha de esa gente de forma virtual a través de los 
    grupos de Telegram. Recomendamos fuertemente estar en esos grupos para enterarse de 
    todas las actividades que se organizan regularmente, están buenísimas.

    \href{https://t.me/comcomdcuba}{Grupo de la ComCom}

    \href{https://t.me/+oDiLv44I7b5jMjZh}{Grupo general del DC}

\end{frame}